% FILE: artifacts/mixture-input-encoding_poster_ccn2026/tex/parts/macros.tex
% ==================================================================================================
% The poster's own layer: the type scale and its math sizes at true scale, the layout grid, the
% full-width region anchors and the helpers that place a built box on the page. The house style
% supplies the type roles, the figure toolkit and the card grammar, and the mixture-input-encoding
% project supplies the notation; parts/preamble.tex loads all of them before this file.

% ===== THE TYPE SYSTEM ============================================================================
% THE ONLY PLACE A POINT SIZE MAY BE CHOSEN. Never write \fontsize elsewhere -- not in a block, not in
% a figure, not inside a TikZ node. If a role needs a size that is not here, add a NAMED macro here.
%
% The poster is 147.32 x 172.72 cm at TRUE SCALE: 1 pt = 0.351 mm, so a point size IS a physical size
% (40 pt = a 14.1 mm em; 24 pt = 8.4 mm). Read every number below as a millimetre statement.
%
% TWO READABILITY FLOORS, read at different distances. Nothing may be set below its own.
%   TEXT   36 pt (12.7 mm em) -- anything read as language: prose, headings, captions, tags, notes,
%                               references, card content, INCLUDING text drawn inside a TikZ card.
%   FIGURE 24 pt ( 8.4 mm em) -- annotations inside a diagram: axis ticks, symbol labels, marks.
%
% The roles that these sizes fill, the two figure levels, the three leadings and the math rule are
% stated once, in artifacts/_common/style/type-roles.tex.
%
% ----- The scale: 8 point sizes ------------------------------------------------------------------
% TEXT (floor 36)
\newcommand{\titlesize}{\fontsize{94}{108}\selectfont}      % poster title
\newcommand{\subtitlesize}{\fontsize{78}{90}\selectfont}    % title continuation
\newcommand{\headingsize}{\fontsize{52}{60}\selectfont}     % block headings, authors, display formulae
\newcommand{\cardtitlesize}{\fontsize{44}{50}\selectfont}   % card banner titles, heading kickers
\newcommand{\bodysize}{\fontsize{40}{56}\selectfont}        % prose that wraps            (PROSE 1.40)
\newcommand{\bodycompact}{\fontsize{40}{46}\selectfont}     % body size, tight            (COMPACT)
\newcommand{\notesize}{\fontsize{36}{42}\selectfont}        % citations, metadata, tags, notes -- FLOOR
% FIGURE (floor 24)
\newcommand{\figcapsize}{\fontsize{36}{42}\selectfont}      % figure captions, shared axis labels, keys
\newcommand{\figannotsize}{\fontsize{30}{35}\selectfont}    % the figure's OBJECTS (vectors, matrices)
\newcommand{\figticksize}{\fontsize{24}{28}\selectfont}     % subordinate marks: ticks, axes -- FLOOR
% ICON GLYPHS -- sized to the disc that holds them, exempt from both floors.
\newcommand{\iconsm}{\fontsize{20}{20}\selectfont}          % 11 mm disc  (\checkring)
\newcommand{\iconmd}{\fontsize{32}{32}\selectfont}          % 28 mm disc  (card banners, bullet marks)
\newcommand{\iconlg}{\fontsize{40}{40}\selectfont}          % 22-24 mm disc (step numbers, question mark)
\newcommand{\iconxl}{\fontsize{48}{48}\selectfont}          % QR legend glyph
\newcommand{\iconxxl}{\fontsize{57}{57}\selectfont}         % block-heading number disc

% --- MATH SIZES, ONE DECLARATION PER STEP OF THE TYPE SCALE ---------------------------------------
% \DeclareMathSizes{<text>}{<math>}{<script>}{<scriptscript>} fires when the TEXT size is selected --
% inside a TikZ node too. EVERY step of the scale above that can carry math needs a line here;
% a step without one falls back to LaTeX's defaults and drifts out of the system SILENTLY.
%
% MATH = 1.125 x TEXT. The body is Source Sans 3, the math Latin Modern (serif): its smaller x-height
% makes a symbol read small beside the sans at the same nominal size, and the 12.5% bump levels them.
%
% SCRIPT / SCRIPTSCRIPT keep the conventional 0.7 / 0.5 -- deliberate, since the floors above
% govern TEXT, not indices. KNOWN CONSEQUENCE: at body size the "mix" index lands at 23 pt, just under
% the 24 pt figure floor and the smallest mark on the poster. If it fails at viewing distance, raise
% the last two numbers here; nothing else changes.
\DeclareMathSizes{24}{27}{19}{14}      % \figticksize  -- figure floor
\DeclareMathSizes{30}{34}{24}{17}      % \figannotsize
\DeclareMathSizes{36}{40}{28}{20}      % \notesize / \figcapsize -- text floor
\DeclareMathSizes{40}{45}{32}{23}      % \bodysize / \bodycompact -- incl. the two move (a) formulae
\DeclareMathSizes{44}{50}{35}{25}      % \cardtitlesize
\DeclareMathSizes{52}{58}{41}{29}      % \headingsize
% The brace readings under the move (a) formulae do NOT follow these sizes -- they are sans, pinned by
% \brlab (artifacts/_common/style/type-roles.tex), because a gloss is words, not algebra.

% ===== Layout grid (cm; origin at the page's south-west corner) ===================================
% Page (true CCN 2026 portrait size).
\def\PW{147.32}\def\PH{172.72}
% Columns (x). The two content columns stack their blocks INDEPENDENTLY (no rigid 2x2),
% which lets each side flow to fit its content.
% THE SIDE MARGIN IS THE ONE CHOSEN NUMBER HERE; the right edge and the full content width are
% ARITHMETIC ON IT AND \PW, not three numbers that have to be kept in step by hand.
\def\gMarg{5}                          % side margin, both sides
\edef\gXl{\gMarg}                      % left content edge
\edef\gXr{\fpeval{\PW-\gMarg}}         % right content edge            (142.32)
\def\gWl{65}           % the width \sbApprBody builds the GLOBAL move boxes at (block-approach.tex)
\edef\gWfull{\fpeval{\PW-2*\gMarg}}    % full content width (footer / supporting)   (137.32)

% ===== REGION MAP (cm; y = NORTH edge of a block, h = its height) =================================
% The two-question layout's own map (\yPhen / \hPhen, \yCol, \gXcolL / \gXcolR / \gGutC) lives in
% poster.tex, because it is specific to that layout. Only the two full-width anchors
% shared with the block files stay here:
\def\yHead{168.72}                           % header banner, full width. Its height is measured, not
                                             % set: read "[layout] header banner height" after any
                                             % header edit and re-derive the column tops.
\def\yFoot{12}                               % footer, full width

% Coordinate helper: \xy{x}{y} -> point at (x cm, y cm) from the page corner (O).
\newcommand{\xy}[2]{($(O)+(#1 cm,#2 cm)$)}

% Drop a pre-built box with its NORTH-WEST corner at (x,y) cm (call INSIDE the picture):
%   \dropbox{<\savebox>}{<x>}{<y>}
\newcommand{\dropbox}[3]{%
  \node[anchor=north west, inner sep=0pt] at \xy{#2}{#3} {\usebox#1};%
}
% As \dropbox but with an explicit anchor (default north west):
%   \dropboxa[<anchor>]{<\savebox>}{<x>}{<y>}
\newcommand{\dropboxa}[4][north west]{%
  \node[anchor=#1, inner sep=0pt] at \xy{#3}{#4} {\usebox#2};%
}

% Shared white space AFTER a body paragraph, before the figure row it introduces: a paragraph must not
% sit tight against the figure it announces.
% The space after a block HEADING is NOT set here -- it is \headbodygap (block-columns.tex), and no
% panel may add to it: heading, then content, nothing in between.
\newcommand{\paragap}{12mm}
